% concatenated from Sunflower.tex
\documentclass[11pt]{article}

\usepackage[a4paper, margin=3.5cm]{geometry}
\usepackage{amsmath, amssymb, amsthm, mathtools}
\usepackage[dvipsnames]{xcolor}
\usepackage{hyperref}
\usepackage{enumitem}
\usepackage{microtype}

\hypersetup{
  colorlinks=true,
  linkcolor=blue!50!black,
  citecolor=blue!50!black,
  urlcolor=blue!50!black
}

\newtheorem{theorem}{Theorem}
\newtheorem{proposition}[theorem]{Proposition}
\newtheorem{lemma}[theorem]{Lemma}
\newtheorem{conjecture}[theorem]{Conjecture}
\theoremstyle{definition}
\newtheorem{remark}[theorem]{Remark}
\newtheorem{example}[theorem]{Example}
\newtheorem{definition}[theorem]{Definition}

\newcommand{\F}{\mathbb{F}}
\newcommand{\rs}{\mathrm{rowspace}\,}
\newcommand{\spn}{\mathrm{span}\,}

\title{On set-like sunflower-free families of subspaces\\over finite fields}
\author{Kamil Otal\\
  \small T\"UB\.ITAK B\.ILGEM UEKAE National Research Institute of Electronics and Cryptology\\
  \small Gebze, Kocaeli, T\"urkiye\\
  \small \texttt{kamil.otal@gmail.com}}

\date{\today}

\begin{document}
\maketitle

\begin{abstract}
The Erd\H{o}s--Rado sunflower problem admits two natural analogues in finite vector spaces, corresponding to two different ways of generalising the set-theoretic notion of a sunflower. The first, used by Ihringer and Kupavskii [FFA 110 (2026) 102746], requires the petals to be in general position over the kernel; the second, used in the subspace codes literature (cf.\ Etzion--Raviv [DAM 186 (2015) 87-97], Blokhuis--De Boeck--D'haeseleer [DCC 90 (2022) 2101-2111]), requires only that the kernel equals the pairwise intersection of distinct petals. We refer to the second version as a \emph{set-like sunflower}, following Ihringer and Kupavskii.

In this note, we focus on the set-like setting. We observe that the constructions of Ihringer--Kupavskii, although correct under their (stronger) definition, do not yield set-like sunflower-free families: we exhibit explicit set-like sunflowers inside their Example~3.1. We then present a construction of set-like $s$-sunflower-free families of $k$-spaces, based on a manipulated version of the lifting construction. To our knowledge, this is the first systematic construction tailored to this setting.

\medskip
\noindent\textbf{MSC:} 05D05, 51E23.

\noindent\textbf{Keywords:} sunflower, $\Delta$-system, subspace code, vector space, $q$-analogue, MRD code, lifting construction.
\end{abstract}

\section{Introduction}

A family of $s$ sets $S_1, \ldots, S_s$ is called an \emph{$s$-sunflower} (or \emph{$\Delta(s)$-system}) with kernel $K$ if $K = S_i \cap S_j$ for all distinct $i, j$.\footnote{The term ``$\Delta$-system'' goes back to Erd\H{o}s and Rado. The term ``sunflower'' appeared in a paper by Deza and Frankl~\cite{deza-frankl} and since the late 1980s has been used by the Boolean circuit complexity community, until it eventually replaced the $\Delta$-system terminology.}\footnote{Here, a family is a set, not a multiset.} A family of $k$-sets is \emph{$s$-sunflower-free} if no $s$ of its members form an $s$-sunflower. In 1960, Erd\H{o}s and Rado~\cite{erdos-rado} proved that for $s \geq 3$ every $s$-sunflower-free family of $k$-sets has size at most $k!(s-1)^k$, and conjectured a much sharper bound:

\begin{conjecture}[Erd\H{o}s--Rado sunflower conjecture]
For every $s \geq 3$ there exists a constant $C$ such that every $s$-sunflower-free family of $k$-sets $\mathcal{F}$ satisfies $|\mathcal{F}| \leq (Cs)^k$.
\end{conjecture}

A complete $k$-partite hypergraph with parts of size $s-1$ gives an $s$-sunflower-free family of $k$-sets of size $(s-1)^k$. A breakthrough of Alweiss, Lovett, Wu, and Zhang~\cite{alwz}, refined slightly by Bell, Chueluecha and Warnke~\cite{bcw}, pushed the bound much closer to the conjectured value:

\begin{proposition}[\cite{alwz,bcw}]
For every $s \geq 3$, any $s$-sunflower-free family of $k$-sets $\mathcal{F}$ satisfies $|\mathcal{F}| \leq (Cs \log k)^k$ for some absolute constant $C$.
\end{proposition}

\subsection{Sunflowers in finite vector spaces}

Let $q$ be a prime power, $\F_q$ the finite field with $q$ elements, and $\F_q^n$ the $n$-dimensional vector space over $\F_q$. When generalising the notion of a sunflower to finite vector spaces, two natural choices arise.

\begin{definition}[General-position sunflower; Ihringer--Kupavskii~\cite{ik2026}]\label{def:gp-sunflower}
A family $S_1, \ldots, S_s$ of $k$-spaces of $\F_q^n$ is a \emph{(general-position) $s$-sunflower} with kernel $K$ if there exists a $d$-space $K$ with $K = S_i \cap S_j$ for all distinct $i, j$, and $S_1/K, \ldots, S_s/K$ are in general position, i.e.\
$\dim(S_1/K + \cdots + S_s/K) = \sum_{i=1}^s \dim(S_i/K)$.
Equivalently, $\dim(S_1 + \cdots + S_s) = d + s(k-d)$.
\end{definition}

\begin{definition}[Set-like sunflower; cf.~Etzion--Raviv~\cite{etzion-raviv}]\label{def:set-like-sunflower}
A family $S_1, \ldots, S_s$ of $k$-spaces of $\F_q^n$ is a \emph{set-like $s$-sunflower} with kernel $K$ if $K = S_i \cap S_j$ for all distinct $i, j$.
\end{definition}

The terminology \emph{set-like sunflower} is due to Ihringer and Kupavskii~\cite{ik2026}, who note that a parallel literature focusing only on pairwise intersections has been developed in the context of subspace codes; see, for example, Etzion and Raviv~\cite{etzion-raviv} and Blokhuis, De Boeck and D'haeseleer~\cite{bdh}.

Every general-position sunflower is a set-like sunflower, but not conversely. The two definitions therefore lead to two distinct extremal problems. In the set case the two notions coincide, because three sets pairwise intersecting in a common set $K$ automatically satisfy the analogous general-position condition; in the vector-space case they do not, due to the failure of the inclusion-exclusion principle for sums of three or more subspaces (see Section~\ref{sec:incl-excl}).

A family of $k$-spaces is \emph{(general-position) $s$-sunflower-free} (respectively, \emph{set-like $s$-sunflower-free}) if it contains no $s$-sunflower in the corresponding sense. Since a set-like sunflower is a more permissive notion, set-like sunflower-freeness is the stronger property: every set-like $s$-sunflower-free family is also general-position $s$-sunflower-free, but not conversely.

\subsection{Related work and contributions}

Ihringer and Kupavskii~\cite{ik2026} studied the Erd\H{o}s--Rado problem under Definition~\ref{def:gp-sunflower}. They adapted the upper bound of Erd\H{o}s and Rado to obtain
\[
|\mathcal{F}| \;\leq\; \prod_{i=1}^k [i(s-1)]_q \;\leq\; \Bigl(\tfrac{q}{q-1}\Bigr)^{\!k} q^{(s-1)\binom{k+1}{2} - k},
\]
and, by an iterative nesting of lifted MRD codes, constructed general-position $s$-sunflower-free families approaching this bound. The upper bound applies to set-like sunflower-free families as well, since the latter form a subclass of the former.

In the subspace codes tradition, the focus has typically been on related notions: equidistant subspace codes (cf.~\cite{etzion-raviv}), subspace packings (cf.~\cite{ekop-subspace-packings}), and sunflower-bound type results for $k$-spaces pairwise intersecting in a point (cf.~\cite{bdh}), all in the spirit of Definition~\ref{def:set-like-sunflower}.

In this note we focus on the set-like setting. Our contributions are:
\begin{itemize}[leftmargin=*, itemsep=2pt]
  \item We show that the constructions of Ihringer--Kupavskii~\cite{ik2026}, although correct as general-position $s$-sunflower-free families, are \emph{not} set-like $s$-sunflower-free; we exhibit explicit set-like $3$-sunflowers in their Example~3.1 (Section~\ref{sec:incl-excl}).
  \item We present a construction of set-like $s$-sunflower-free families of $k$-spaces of $\F_q^n$ with size $q^{n-k}$, based on a manipulated version of the lifting construction (Section~\ref{sec:construction}). To our knowledge this is the first systematic construction tailored to this setting.
\end{itemize}

\section{Preliminaries}

\subsection{Maximum rank distance codes}

Let $\F_q^{k \times \ell}$ denote the set of $k \times \ell$ matrices over $\F_q$. The \emph{rank distance} on $\F_q^{k \times \ell}$ is $d_R(A, B) := \mathrm{rank}(A - B)$. A subset $\mathcal{C} \subseteq \F_q^{k \times \ell}$ with $|\mathcal{C}| \geq 2$ is a \emph{rank-metric code}; its \emph{minimum rank distance} is $d_R(\mathcal{C}) := \min\{d_R(A,B) : A, B \in \mathcal{C}, A \neq B\}$.

\begin{proposition}[Singleton-like bound~\cite{delsarte}]
For any rank-metric code $\mathcal{C} \subseteq \F_q^{k \times \ell}$,
\[
|\mathcal{C}| \;\leq\; q^{\max\{k, \ell\}(\min\{k, \ell\} - d_R(\mathcal{C}) + 1)}.
\]
\end{proposition}

A rank-metric code attaining this bound is called a \emph{maximum rank distance (MRD) code}.

One useful family of MRD codes comes from the matrix representation of $\F_{q^k}$. Let $f$ be a monic irreducible polynomial of degree $k$ over $\F_q$, and $A$ its $k \times k$ companion matrix. Then $\spn(I, A, \ldots, A^{k-1})$, the set of $\F_q$-polynomials in $A$ of degree less than $k$, is a faithful matrix representation of $\F_{q^k}$; see~\cite[\S2.5]{lidl}.

\begin{proposition}\label{prop:matrix-rep}
Let $\mathcal{C} \subseteq \F_q^{k \times k}$ be the matrix representation of $\F_{q^k}$. Then $\mathcal{C}$ is an MRD code with $d_R(\mathcal{C}) = k$.
\end{proposition}

The proof is immediate: $\mathcal{C}$ is linear, and every nonzero element is invertible.

\subsection{Subspace codes and lifting}

Let $\mathcal{P}_q(n)$ denote the set of all subspaces of $\F_q^n$. The \emph{subspace distance} is
\[
d_S(U, V) := \dim U + \dim V - 2 \dim(U \cap V).
\]
A subset $\mathcal{F} \subseteq \mathcal{P}_q(n)$ with $|\mathcal{F}| \geq 2$ is a \emph{subspace code}. If all members have the same dimension $k$, then $\mathcal{F} \subseteq \mathcal{G}_q(n, k)$ is a \emph{constant-dimension code}.

The \emph{lifting construction} produces large constant-dimension codes from MRD codes:

\begin{proposition}[Silva--Kschischang--K\"otter~\cite{silva}]\label{prop:lifting}
Let $\mathcal{C} \subseteq \F_q^{k \times (n-k)}$ be an MRD code with $d_R(\mathcal{C}) = \min\{k, n-k\}$, and let $I$ denote the $k \times k$ identity matrix. Then
\[
\mathcal{F}_{\mathcal{C}} \;=\; \{ \rs [I \mid A] : A \in \mathcal{C}\}
\]
is a constant-dimension subspace code of size $q^{\max\{k, n-k\}(\min\{k, n-k\} - d_R(\mathcal{C}) + 1)}$ and minimum distance $d_S(\mathcal{F}_{\mathcal{C}}) = 2 d_R(\mathcal{C})$.
\end{proposition}

\section{The role of inclusion-exclusion}\label{sec:incl-excl}

For two subspaces $U, V$ of a finite-dimensional vector space,
\[
\dim(U + V) = \dim U + \dim V - \dim(U \cap V),
\]
mirroring the inclusion-exclusion identity for two sets. The analogue \emph{fails} for three subspaces: in general,
\begin{align*}
\dim(U + V + W) \;\neq\; & \dim U + \dim V + \dim W \\
& - \dim(U \cap V) - \dim(U \cap W) - \dim(V \cap W) + \dim(U \cap V \cap W).
\end{align*}
For example, take $U = \spn\{(1,0,0), (0,1,0)\}$, $V = \spn\{(1,0,0), (0,0,1)\}$, $W = \spn\{(1,0,0), (0,1,1)\}$ in $\F_q^3$. Each pair intersects in $\spn\{(1,0,0)\}$, the triple intersection is the same line, yet $U + V + W = \F_q^3$, so $\dim(U+V+W) = 3$, whereas the right-hand side above evaluates to $2 + 2 + 2 - 1 - 1 - 1 + 1 = 4$.

This failure is precisely what separates Definitions~\ref{def:gp-sunflower} and~\ref{def:set-like-sunflower}. In a general-position $s$-sunflower the identity $\dim(S_1+\cdots+S_s) = d + s(k-d)$ holds by definition (the general-position condition is exactly this equality). In a set-like sunflower no such identity is required, and the sum dimension may be strictly less than $d + s(k-d)$.

\subsection{A set-like sunflower inside the Ihringer--Kupavskii construction}

We illustrate the gap between the two definitions by exhibiting a set-like sunflower inside the smallest explicit construction of~\cite{ik2026}.

\begin{example}[\cite{ik2026}, Example 3.1]\label{ex:ik-example}
Fix a $1$-space $T$ in $\F_q^5$, and choose $q^2 + 1$ $3$-spaces $\Pi_1, \ldots, \Pi_{q^2+1}$ through $T$ with $\Pi_i \cap \Pi_j = T$ for $i \neq j$. In $\Pi_1$ take all $q^2 + q + 1$ $2$-spaces; in each $\Pi_i$ with $i \geq 2$ take the $q^2$ $2$-spaces disjoint from $T$. The resulting family $\mathcal{F}$ has $q^4 + q^2 + q + 1$ members.
\end{example}

Ihringer and Kupavskii prove in~\cite{ik2026} that $\mathcal{F}$ is general-position $3$-sunflower-free. Their argument crucially uses the equality $\dim(S_1+S_2+S_3) = d + 3(k-d)$, which holds for general-position sunflowers by definition. The next example shows that $\mathcal{F}$ contains set-like $3$-sunflowers, so the same family is \emph{not} set-like $3$-sunflower-free.

\begin{example}\label{ex:counter}
Take $q = 2$, $T = \spn\{[1,0,0,0,0]\}$, and
\[
\Pi_1 = \rs\begin{pmatrix} 1 & 0 & 0 & 0 & 0 \\ 0 & 0 & 1 & 1 & 0 \\ 0 & 0 & 0 & 0 & 1 \end{pmatrix}.
\]
The three $2$-spaces of $\Pi_1$ that contain $T$ are
\[
S_1 = \rs\begin{pmatrix} 1 & 0 & 0 & 0 & 0 \\ 0 & 0 & 1 & 1 & 0 \end{pmatrix}, \quad
S_2 = \rs\begin{pmatrix} 1 & 0 & 0 & 0 & 0 \\ 0 & 0 & 0 & 0 & 1 \end{pmatrix},
\]
\[
S_3 = \rs\begin{pmatrix} 1 & 0 & 0 & 0 & 0 \\ 0 & 0 & 1 & 1 & 1 \end{pmatrix}.
\]
A direct check shows $S_i \cap S_j = T$ for all distinct $i, j$, so $\{S_1, S_2, S_3\}$ is a set-like $3$-sunflower with kernel $T$. Note that $\dim(S_1 + S_2 + S_3) = \dim(\Pi_1) = 3 < 1 + 3 \cdot 1 = 4$, so $\{S_1, S_2, S_3\}$ is not a general-position sunflower, consistent with the analysis of~\cite{ik2026}.
\end{example}

More generally, the same phenomenon occurs for arbitrary $q \geq 2$: the $q+1$ $2$-spaces of $\Pi_1$ that contain $T$ form a pencil, and any three of them constitute a set-like $3$-sunflower with kernel $T$. Analogous set-like sunflowers occur in the more involved constructions of~\cite[\S3.3, \S3.4]{ik2026}.

This is not a contradiction with~\cite{ik2026}: their constructions are designed under Definition~\ref{def:gp-sunflower} and remain correct under that definition. The observation is rather that the two extremal problems are genuinely distinct, and that an alternative construction is required for the set-like setting.

\section{A construction for set-like sunflower-free families}\label{sec:construction}

We now give a construction of set-like $s$-sunflower-free families of $k$-spaces. The construction is inspired by the construction of almost affinely disjoint (AAD) spaces in~\cite{aad-isit, aad-2024, aad-ffa}; in particular, the simultaneous use of $A$ and $A^2$ originates from these references.

We use the following standard identity:

\begin{proposition}\label{prop:sum-complement}
Let $U_1, U_2, \ldots, U_s$ be subspaces of $\F_q^n$. Then
\[
(U_1 + U_2 + \cdots + U_s)^\perp \;=\; U_1^\perp \cap U_2^\perp \cap \cdots \cap U_s^\perp.
\]
\end{proposition}

For a proof see, e.g.,~\cite{otal-thesis}. This identity is a standard tool in the subspace codes literature, used for instance in~\cite{ekop-subspace-packings}.

\begin{theorem}\label{thm:construction}
Let $n \geq 2\ell + 1$ and $k = n - \ell$. Let $I$ denote the $\ell \times \ell$ identity matrix over $\F_q$, and let $\mathcal{C} \subseteq \F_q^{\ell \times \ell}$ be the matrix representation of $\F_{q^\ell}$ over $\F_q$, as in Proposition~\ref{prop:matrix-rep}. For a matrix $A$, write $[A]_1$ for its first column, and let $\mathbf{0}$ denote the $\ell \times (n - 2\ell - 1)$ zero matrix. Define
\[
\mathcal{G} \;=\; \bigl\{ \rs [I \mid A \mid [A^2]_1 \mid \mathbf{0}] : A \in \mathcal{C}\bigr\}.
\]
Then $\mathcal{F} = \{U^\perp : U \in \mathcal{G}\}$ is a set-like $s$-sunflower-free family of $k$-spaces of $\F_q^n$ of size $q^\ell = q^{n-k}$ for every $s \geq 3$.
\end{theorem}

\begin{proof}
Since $\mathcal{C}$ has $q^\ell$ elements, $|\mathcal{F}| = q^\ell$. Each $U \in \mathcal{G}$ is an $\ell$-space, so each $U^\perp$ is a $(n-\ell) = k$-space.

A set-like $3$-sunflower-free family is automatically set-like $s$-sunflower-free for $s \geq 3$ (any $s$-sunflower contains a $3$-sub-sunflower with the same kernel). It therefore suffices to show that $\mathcal{F}$ is set-like $3$-sunflower-free.

By Proposition~\ref{prop:sum-complement}, for any $U, V, W \in \mathcal{F}$ corresponding to $A, B, C \in \mathcal{C}$,
\[
\dim(U \cap V) = n - \dim(U^\perp + V^\perp), \qquad \dim(U \cap V \cap W) = n - \dim(U^\perp + V^\perp + W^\perp).
\]

\textit{Pairwise intersection.} For distinct $A, B \in \mathcal{C}$, row reduction gives
\[
\mathrm{rank}\begin{pmatrix} I & A & [A^2]_1 & \mathbf{0} \\ I & B & [B^2]_1 & \mathbf{0} \end{pmatrix}
= \mathrm{rank}\begin{pmatrix} I & A & [A^2]_1 & \mathbf{0} \\ 0 & B - A & [B^2 - A^2]_1 & \mathbf{0} \end{pmatrix}
= 2\ell,
\]
since $B - A$ is invertible (Proposition~\ref{prop:matrix-rep}). Hence $\dim(U \cap V) = n - 2\ell$.

\textit{Triple intersection.} For pairwise distinct $A, B, C \in \mathcal{C}$,
\begin{align*}
\mathrm{rank}\begin{pmatrix} I & A & [A^2]_1 & \mathbf{0} \\ I & B & [B^2]_1 & \mathbf{0} \\ I & C & [C^2]_1 & \mathbf{0} \end{pmatrix}
&\;=\; \mathrm{rank}\begin{pmatrix} I & A & [A^2]_1 & \mathbf{0} \\ 0 & B - A & [B^2 - A^2]_1 & \mathbf{0} \\ 0 & C - A & [C^2 - A^2]_1 & \mathbf{0} \end{pmatrix} \\
&\;=\; \mathrm{rank}\begin{pmatrix} I & A & [A^2]_1 & \mathbf{0} \\ 0 & I & [B + A]_1 & \mathbf{0} \\ 0 & I & [C + A]_1 & \mathbf{0} \end{pmatrix} \\
&\;=\; \mathrm{rank}\begin{pmatrix} I & A & [A^2]_1 & \mathbf{0} \\ 0 & I & [B + A]_1 & \mathbf{0} \\ 0 & 0 & [C - B]_1 & \mathbf{0} \end{pmatrix} \;=\; 2\ell + 1.
\end{align*}
Here we used that $\mathcal{C}$ is a field (so closed under multiplication, inverses, and so that $(B-A)^{-1}(B^2 - A^2) = B + A$), together with the column-extraction identities $X[Y]_1 = [XY]_1$ and $[X]_1 + [Y]_1 = [X+Y]_1$ valid for all $X, Y \in \mathcal{C}$. Since $C \neq B$ implies $C - B$ is invertible, $[C-B]_1$ is a nonzero column and contributes rank $1$. Hence $\dim(U \cap V \cap W) = n - (2\ell + 1) = n - 2\ell - 1$.

\textit{Conclusion.} For pairwise distinct $U, V, W \in \mathcal{F}$,
\[
\dim(U \cap V \cap W) \;=\; n - 2\ell - 1 \;<\; n - 2\ell \;=\; \dim(U \cap V).
\]
If $\{U, V, W\}$ were a set-like $3$-sunflower with kernel $K$, then $K = U \cap V = U \cap W = V \cap W$, which forces $K = U \cap V \cap W$ and hence $\dim(U \cap V) = \dim(U \cap V \cap W)$, contradicting the strict inequality above. Therefore $\mathcal{F}$ contains no set-like $3$-sunflower.
\end{proof}

\begin{remark}
The smallest meaningful instance of Theorem~\ref{thm:construction} has parameters $q = 2$, $n = 5$, $\ell = 2$ (so $k = 3$), giving a set-like $s$-sunflower-free family of $4$ three-spaces in $\F_2^5$.
\end{remark}

\section{Concluding remarks and open problems}

The construction in Theorem~\ref{thm:construction} produces set-like $s$-sunflower-free families of size $q^{n-k}$. This is far below the upper bound $\prod_{i=1}^k [i(s-1)]_q$ inherited from~\cite{ik2026}, which applies to the set-like setting as well: the proof of that upper bound (a Erd\H{o}s--Rado-style pigeonhole over $1$-subspaces) uses only the pairwise intersection structure of a sunflower and therefore goes through verbatim under Definition~\ref{def:set-like-sunflower}.

Several natural questions remain:

\begin{enumerate}[leftmargin=*, itemsep=2pt]
  \item \textbf{Order of magnitude.} What is the order of magnitude of the largest set-like $s$-sunflower-free family of $k$-spaces of $\F_q^n$? Is it polynomially, or exponentially, smaller than the largest general-position $s$-sunflower-free family?
  \item \textbf{Iterated constructions.} Can the construction of Theorem~\ref{thm:construction} be iterated or nested, in the spirit of the iterative MRD-nesting of~\cite{ik2026}, in a way that preserves set-like sunflower-freeness?
  \item \textbf{Dependence on $n$.} The set-like problem is genuinely dimension-dependent (as is also the case for the related sunflower-bound problems studied in~\cite{bdh}). For fixed $k$ and $s$, what is the dependence of the extremal size on $n$?
  \item \textbf{Bridging the two definitions.} Is there a natural intermediate notion (for instance, requiring general position only up to triples, or only on the level of $r$-fold sums for fixed $r$) that interpolates between the two regimes?
\end{enumerate}

%\section*{Acknowledgements}

%The author thanks Ferdinand Ihringer and Andrey Kupavskii for their stimulating paper~\cite{ik2026}, which prompted the present investigation.

\begin{thebibliography}{99}

\bibitem{alwz} R.~Alweiss, S.~Lovett, K.~Wu, J.~Zhang, Improved bounds for the sunflower lemma, \emph{Ann. of Math.} (2) \textbf{194} (2021), 795--815.

\bibitem{aad-isit} T.~Arikan, B.~D\"uzg\"un, K.~Otal, F.~\"Ozbudak, A new construction of asymptotically optimal almost affinely disjoint spaces, in \emph{IEEE ISIT 2023}, pp.~282--285.

\bibitem{bcw} T.~Bell, S.~Chueluecha, L.~Warnke, Note on sunflowers, \emph{Discrete Math.} \textbf{344} (2021), no.~7.

\bibitem{bdh} A.~Blokhuis, M.~De Boeck, J.~D'haeseleer, On the sunflower bound for $k$-spaces, pairwise intersecting in a point, \emph{Des. Codes Cryptogr.} \textbf{90} (2022), 2101--2111.

\bibitem{delsarte} P.~Delsarte, Bilinear forms over a finite field, with applications to coding theory, \emph{J. Combin. Theory Ser. A} \textbf{25} (1978), 226--241.

\bibitem{deza-frankl} M.~Deza, P.~Frankl, Every large set of equidistant $(0, +1, -1)$-vectors forms a sunflower, \emph{Combinatorica} \textbf{1} (1981), 225--231.

\bibitem{aad-2024} B.~D\"uzg\"un, T.~Arikan, K.~Otal, F.~\"Ozbudak, Asymptotically optimal $[2k+1, k, k]$-almost affinely disjoint subspaces, \emph{Discrete Math.} \textbf{347} (2024), 114140.

\bibitem{erdos-rado} P.~Erd\H{o}s, R.~Rado, Intersection theorems for systems of sets, \emph{J. London Math. Soc.} (2) \textbf{35} (1960), 85--90.

\bibitem{ekop-subspace-packings} T.~Etzion, S.~Kurz, K.~Otal, F.~\"Ozbudak, Subspace packings: constructions and bounds, \emph{Des. Codes Cryptogr.} \textbf{88} (2020), 1781--1810.

\bibitem{etzion-raviv} T.~Etzion, N.~Raviv, Equidistant codes in the Grassmannian, \emph{Discrete Appl. Math.} \textbf{186} (2015), 87--97.

\bibitem{ik2026} F.~Ihringer, A.~Kupavskii, The Erd\H{o}s--Rado sunflower problem for vector spaces, \emph{Finite Fields Appl.} \textbf{110} (2026), 102746.

\bibitem{lidl} R.~Lidl, H.~Niederreiter, \emph{Finite Fields}, Encyclopedia of Mathematics and its Applications \textbf{20}, 2nd ed., Cambridge University Press, 1997.

\bibitem{otal-thesis} K.~Otal, \emph{Constructions of Maximum Rank Distance Codes, Cyclic Constant Dimension Codes, and Subspace Packings}, PhD thesis, Middle East Technical University, 2018.

\bibitem{aad-ffa} K.~Otal, T.~Arikan, An asymptotically optimal construction of almost affinely disjoint subspaces, \emph{Finite Fields Appl.} \textbf{84} (2022), 102099.

\bibitem{silva} D.~Silva, F.~Kschischang, R.~K\"otter, A rank-metric approach to error control in random network coding, \emph{IEEE Trans. Inform. Theory} \textbf{54} (2008), 3951--3967.

\end{thebibliography}

\end{document}
